\uuid{koDg}
\titre{Étude de fonction}
\auteur{bodin ; hanani ; mzari}
\organisation{exo7}
\licence{CC BY-NC-SA 4.0}
\datecreate{2020-01-20}
\niveau{L1}
\module{Analyse}
\chapitre{Fonctions circulaires et hyperboliques inverses}
\sousChapitre{Autre}
\difficulte{2}
\theme{Fonctions usuelles, étude de fonction}
\exo7qcmid{lille-1.fonctions-usuelles.015}

\contenu{
\begin{question}
\qclassification{126.00}
\qtags{motcle=étude de fonction}
Soit $f(x)=\arcsin x + \arccos x$.  Quelles sont les assertions vraies ?

\begin{answers}
    \good{Le domaine de définition de $f$ est $[-1,1]$.}

    \good{$\forall x\in [-1,1], \, f(x)=\frac{\pi}{2}$}

    \bad{$\forall x\in [-1,1], \, f(x)= x$}

    \good{$f$ est une fonction constante.}
\end{answers}
\begin{explanations}
$f$ est définie sur $[-1,1]$,  dérivable sur $]-1,1[$  et $f'(x)=0$, $\forall x \in ]-1,1[$. Puisque $]-1,1[$ est un intervalle, on déduit que $f$ est constante sur $]-1,1[$ et comme $f$ est continue sur $[-1,1]$, $f$ est constante sur $[-1,1]$. Or $f(0)=\frac{\pi}{2}$, donc $f(x)=\frac{\pi}{2}, \, \forall x \in [-1,1]$.
\end{explanations}
\end{question}
}
